\documentclass{article}
\usepackage{amsmath,amssymb}
\usepackage{CJKutf8}
\usepackage[a4paper,margin=22mm]{geometry}

\setlength{\parindent}{0pt}
\setlength{\parskip}{0.55em}

\newcommand{\mondai}[1]{%
  \par\medskip
  \noindent{\large #1}\quad\ignorespaces
}

\newcommand{\defitem}[1]{%
  \par\noindent #1\par
}

\begin{document}
\begin{CJK}{UTF8}{min}

\begin{center}
{\Large 第4章 演習問題}\\
問31(1)(4)(6), 問32(2)(6), 問34(3), 問35(3)
\end{center}

\fbox{\Large 定義}



\fbox{\Large 問題}
問1

次の命題が成り立つか判定せよ。理由も述べよ（証明でなくてよい）。
\begin{enumerate}
  \item $\forall x\in\mathbb{Z}\ \forall y\in\mathbb{N},\ x\geq -y*y$
  \item $\forall x\in\mathbb{Z}\ \exists y\in\mathbb{Z},\ x=y*y$
\end{enumerate}

問2
集合の規則に従って $A\cap (B \cup C) = (A \cap B) \cup (A \cap C)$を示せ。つまり
$$ x \in A \cap (B \cup C) \Rightarrow x \in (A \cap B) \cup (A \cap C) $$
$$ x \in (A \cap B) \cup (A \cap C) \Rightarrow x \in A \cap (B \cup C) $$
をそれぞれ示せ. 
問3
\begin{enumerate}
\item  f, g が全単射であれば、合成 g ◦ f も全単射であることを証明せよ。
\item 合成 g ◦ f が全射であっても、f が全射となるとは限らないことを示せ。
 \end{enumerate}
問4有界な数列 (an)
∞
n=1 に対して An := {am | m ≥ n} とおき、sup
m≥n
an := sup An,
この時, 数列 (an)
∞
n=1 の上極限 lim sup an := limn→∞
sup
m≥n
anが存在することを示せ。
問5  命題 4.12 を証明せよ。すなわち、アルキメデスの公理と実数の完備性から、実
数の連続性の公理を導け。
\end{CJK}\end{document}
